% ==============================================================================
% Περίληψη (Greek Abstract)
% ==============================================================================

\chapter*{Περίληψη}
\addcontentsline{toc}{chapter}{Περίληψη}

Η παρούσα διπλωματική εργασία πραγματεύεται τον σχεδιασμό, την υλοποίηση και την πειραματική αξιολόγηση ενός πλήρως κατανεμημένου συστήματος γενετικών αλγορίθμων (Distributed Genetic Algorithm - DGA) για τη βελτιστοποίηση του αεροδυναμικού σχεδιασμού 3D-εκτυπώσιμων μη επανδρωμένων αεροσκαφών (UAV). Η υπολογιστική αξιολόγηση της αεροδυναμικής συμπεριφοράς μέσω κλασικών μεθόδων προσομοίωσης αποτελεί διαχρονικά το κρισιμότερο υπολογιστικό εμπόδιο, καθώς απαιτεί χιλιάδες χρονοβόρες αξιολογήσεις.

Προκειμένου να επιτευχθεί δραστική επιτάχυνση χωρίς απώλεια της βέλτιστης λύσης, προτείνεται μια πρωτότυπη ιεραρχική αρχιτεκτονική αξιολόγησης τριών επιπέδων (Multi-Tier $\epsilon$-Bypass Pipeline). Στο πρώτο επίπεδο (Tier 1), ανιχνεύονται σχεδόν ταυτόσημες γεωμετρίες με προηγούμενα αξιολογημένα σχέδια και ανακτάται άμεσα η τιμή καταλληλότητας με μηδενικό υπολογιστικό κόστος. Στο δεύτερο επίπεδο (Tier 2), αξιοποιείται ένα ενεργό προσεγγιστικό μοντέλο (Surrogate Model) βασισμένο σε πολυεπίπεδα νευρωνικά δίκτυα (MLP) και διαδικασίες Gauss (Gaussian Processes) με επιγραμμική επανεκπαίδευση (online retraining). Στο τρίτο επίπεδο (Tier 3), αποστέλλονται για πλήρη προσομοίωση μέσω VLM (Vortex Lattice Method) σε εξειδικευμένους Python workers μόνο τα πραγματικά ανεξερεύνητα υποψήφια σχέδια.

Η χωρική οργάνωση και η αποδοτική αναζήτηση γειτονικών σχεδίων στον πολυδιάστατο παραμετρικό χώρο υποστηρίζεται από έναν κατανεμημένο πίνακα κατακερματισμού LEAD DHT (Learned Exact And Approximate DHT). Το σύστημα χαρτογραφεί τον 10-διάστατο χώρο παραμέτρων σε μία διάσταση μέσω πολυσημειακών καμπυλών Hilbert (multi-probe Hilbert curves) και εφαρμόζει μοντέλα ευρετηρίασης RMI (Recursive Model Indexes) με προσαρμοστικό ρυθμιστή PID. Το σύστημα οργανώνεται σε αρχιτεκτονική μικροϋπηρεσιών υψηλής απόδοσης γραμμένο εξ ολοκλήρου σε Rust, διασυνδεδεμένο με διακομιστή εξισορρόπησης φορτίου Envoy gRPC και υποδομή παρακολούθησης ροών Kafka/Fluent-Bit. Τα πειραματικά αποτελέσματα καταδεικνύουν επιτάχυνση του εξελικτικού κύκλου κατά τάξεις μεγέθους, διατηρώντας ταυτόχρονα υψηλή αεροδυναμική απόδοση και ευστάθεια.

\vspace{1.2cm}

\noindent \textbf{Λέξεις-κλειδιά:} Γενετικοί Αλγόριθμοι, Μοντέλο Νήσων, Κατανεμημένα Συστήματα, LEAD DHT, Ευρετήρια Μάθησης, Προσεγγιστικά Μοντέλα (Surrogate Models), Αεροδυναμική Βελτιστοποίηση, AeroSandbox, Rust, gRPC.

\clearpage
